\begin{mistake}{2026-04-09}{34}

\begin{problemcontent}
求 \( f(x) = e^x \sin x \) 的 \( n \) 阶导数 \( f^{(n)}(x) \)。
\end{problemcontent}

\begin{solutioncontent}
答案：\( f^{(n)}(x) = (\sqrt{2})^n e^x \sin\left(x + \frac{n\pi}{4}\right) \)

推导过程：
先求一阶导数，运用乘法求导法则和辅助角公式化简：
\[
\begin{aligned}
f'(x) &= (e^x)'\sin x + e^x(\sin x)' \\
&= e^x\sin x + e^x\cos x \\
&= e^x(\sin x + \cos x) \\
&= \sqrt{2}e^x\sin\left(x + \frac{\pi}{4}\right)
\end{aligned}
\]
可以发现，每次对函数求导，都会在原函数前面提取出系数 \( \sqrt{2} \)，并使正弦项的相位增加 \( \frac{\pi}{4} \)，而函数的整体形式"\( \text{常数} \cdot e^x \cdot \sin(\dots) \)"保持不变。
由此运用数学归纳法可直接得出 \( n \) 阶导数的通项：
\[
f^{(n)}(x) = (\sqrt{2})^n e^x \sin\left(x + \frac{n\pi}{4}\right)
\]
\end{solutioncontent}

\mistakeknowledge{
\begin{itemize}
 \item \textbf{核心公式}：三角函数辅助角公式 \( \sin x + \cos x = \sqrt{2}\sin(x+\frac{\pi}{4}) \)。
 \item \textbf{易错点}：若生搬硬套莱布尼茨公式展开，结果会得到求和级数形式，且在考试有限时间内极难化简为闭合的通项公式。
 \item \textbf{关联知识}：复数法本质。\( e^x\sin x \) 是 \( e^{(1+i)x} \) 的虚部，每次求导乘 \( (1+i) = \sqrt{2}e^{i\pi/4} \)，即模长乘 \( \sqrt{2} \)，辐角加 \( \pi/4 \)。
\end{itemize}}

\end{mistake}
